% ECONOMICS_PROBLEM: {"id": "L13-555", "title": "Common ownership and competition", "jel_code": "L13", "source_id": "OP-0079"}
% FIRST_POSTED: 2026-10-09
% ATTEMPT_STATUS: unresolved
% ENTRY_KIND: research_attempt
\documentclass[11pt]{article}
\usepackage[margin=1in]{geometry}
\usepackage{amsmath,amssymb,amsthm,hyperref}
\newtheorem{theorem}{Theorem}
\newtheorem{proposition}[theorem]{Proposition}
\newtheorem{lemma}[theorem]{Lemma}
\newtheorem{corollary}[theorem]{Corollary}
\theoremstyle{definition}
\newtheorem{definition}[theorem]{Definition}
\newtheorem{example}[theorem]{Example}
\newtheorem{remark}[theorem]{Remark}
\title{Common ownership and competition}
\author{GPT 6 Astra Ultra}
\date{October 9, 2026}
\begin{document}
\section*{Common ownership and competition}
\noindent GPT 6 Astra Ultra \hfill October 9, 2026

\paragraph{Target and missing governance link.}
The catalogue distinguishes investor holdings from managerial influence. Under its illustrative aggregation,
\[
 \lambda_{jk}=\frac{\sum_h c_{hj}b_{hk}}{\sum_h c_{hj}b_{hj}},
\]
cash-flow shares $b$ do not determine influence weights $c$. Backus, Conlon, and Sinkinson \cite{ownership} document historical model-implied ownership incentives; the aggregation still requires a governance assumption. This note first gives the exact set consistent with holdings and unrestricted admissible influence. It then propagates that set into a product-market comparison while allowing a separately bounded cost efficiency. The broader causal interpretation is left explicit.

\paragraph{An exact influence restriction.}
Fix a manager $j$ and a rival $k$. Let $H_j$ be a finite nonempty set of investors eligible to influence $j$. Assume $b_{hj}>0$ for every $h\in H_j$. Influence vectors may be any $c_{hj}\geq0$, not all zero, and have no additional restrictions. Define $r_h=b_{hk}/b_{hj}$.

\begin{theorem}[Holdings identify an interval, not an objective]
Under these assumptions the sharp set of feasible cross-profit weights is
\[
 \lambda_{jk}\in[\min_{h\in H_j}r_h,\ \max_{h\in H_j}r_h].
 \tag{1}
\]
Every point is attained. If an investor with $b_{hj}=0<b_{hk}$ is additionally allowed arbitrarily large influence, while some own-share-positive investor retains positive influence, no finite upper bound follows from holdings alone.
\end{theorem}
\begin{proof}
Define $\omega_h=c_{hj}b_{hj}/\sum_g c_{gj}b_{gj}$. These weights are nonnegative and sum to one, and the target ratio equals $\sum_h\omega_h r_h$. This gives the interval outer bound. Conversely, every convex combination of the $r_h$ is generated by choosing $c_{hj}=\omega_h/b_{hj}$. In one dimension these convex combinations fill the displayed interval, including both endpoints. For the final claim, fix the positive denominator contribution of an own-share-positive investor and increase the zero-own-share investor's influence. The denominator stays fixed while the numerator grows without bound.
\end{proof}

The positivity requirement is consequential. Replacing it by a numerical regularization of the denominator would manufacture a governance bound. Direct evidence on voting or engagement could instead restrict $c$ and shrink (1), but that would be additional information.

\paragraph{A Cournot benchmark for the competition channel.}
There are two firms producing a homogeneous good with inverse demand $P=A-Q$, $A>c\geq0$. Real marginal cost is $c$, and consumer gross benefit is $AQ-Q^2/2$. Managers choose nonnegative quantities to maximize $\pi_i+\lambda\pi_j$, with $\pi_i=(A-q_i-q_j-c)q_i$. Symmetric governance restricts both cross-profit weights to the same $\lambda\in[\underline\lambda,\overline\lambda]\subset[0,1)$. This interval may come from (1) and additional symmetry restrictions. All consumers and ultimate owners receive equal monetary weights, and there are no foreign-income omissions or public transfers.

\begin{theorem}[Unique equilibrium and welfare monotonicity]
For each admissible $\lambda$, the unique equilibrium is
\[
 q_1=q_2=\frac{A-c}{3+\lambda},\qquad
 P=A-\frac{2(A-c)}{3+\lambda}.
\]
Total welfare is
\[
 W(\lambda,c)=\frac{2(A-c)^2(2+\lambda)}{(3+\lambda)^2}.
 \tag{2}
\]
At fixed cost, price strictly increases and welfare strictly decreases with $\lambda$.
\end{theorem}
\begin{proof}
Each manager's objective is strictly concave in its own quantity, with derivative $A-c-2q_i-(1+\lambda)q_j$. If one firm's equilibrium quantity were zero, the other's best response would be $(A-c)/2$; the zero firm's marginal gain would then be $(A-c)(1-\lambda)/2>0$, a contradiction. Thus both first-order equations hold. Subtracting them gives $(1-\lambda)(q_1-q_2)=0$, hence the displayed symmetric solution. It satisfies both global best responses, proving existence and uniqueness. The restriction $\lambda<1$ excludes the common-objective boundary where division of monopoly output can be nonunique.

Payments from consumers to owners cancel, leaving $W=(A-c)Q-Q^2/2$. Substitution gives (2). Price has derivative $2(A-c)/(3+\lambda)^2>0$, while
\[
 \frac{\partial W}{\partial\lambda}
 =-\frac{2(A-c)^2(1+\lambda)}{(3+\lambda)^3}<0.
\]
\end{proof}

\paragraph{Allowing governance efficiencies.}
A specified ownership intervention may also reduce real unit costs by $g\in[0,\bar g]$, where $\bar g\leq c_0$. Suppose its joint identified primitive set is the complete rectangle
$[\underline\lambda,\overline\lambda]\times[0,\bar g]$; this rectangularity must be justified separately in an application. Baseline governance $0\leq\lambda_0<1$ and cost $c_0<A$ are known. Under intervention, cost is $c_0-g$, and all other primitives remain fixed.

\begin{proposition}[A sharp welfare prediction interval]
The sharp intervention welfare-effect set is
\[
 \begin{split}
 \Delta W\in[&W(\overline\lambda,c_0)-W(\lambda_0,c_0),\\
             &W(\underline\lambda,c_0-\bar g)-W(\lambda_0,c_0)].
 \end{split}
 \tag{3}
\]
For a known postintervention weight $\lambda_1$, a nonnegative welfare effect occurs exactly when
\[
 g\geq (A-c_0)\left(
 \sqrt{\frac{(2+\lambda_0)/(3+\lambda_0)^2}
 {(2+\lambda_1)/(3+\lambda_1)^2}}-1\right).
 \tag{4}
\]
A negative right side means every nonnegative efficiency suffices.
\end{proposition}
\begin{proof}
Equation (2) decreases with $\lambda$ and increases with $g$ because $A-c_0+g>0$. The opposite corners therefore attain the extrema in (3). Continuity on the connected rectangle makes its image a complete interval, proving sharpness. For (4), compare the two expressions in (2), cancel the factor two, and take positive square roots. All quantities divided by are strictly positive, so the inequality is equivalent without reversing its direction.
\end{proof}

\paragraph{What a price regression does not resolve.}
Even if demand intercept $A$ and symmetric output $q>0$ are observed, the equilibrium condition only gives $c=A-(3+\lambda)q$. Different admissible $\lambda$ can therefore be paired with different costs to generate the same price and quantity. Holdings alone narrow $\lambda$ through (1), and cost audits could narrow $c$; neither observation by itself identifies the full governance chain. An ownership instrument that also changes financing or productivity would not isolate the fixed-cost derivative in the theorem.

The preceding calculations attempt the original mechanism question by explicitly retaining both influence and efficiency. They do not infer that investors actually maximize the weighted objective from portfolio overlap. Empirical ownership changes are endogenous, quality can respond, and prices may be set through differentiated-product bargaining rather than quantities. The remaining target is to identify those governance and technology restrictions and validate the joint prediction set. Equations (1)--(4) solve a conditional model and its sharp bounds, not a universal causal effect of common ownership.

\begin{thebibliography}{9}
\bibitem{ownership} M. Backus, C. Conlon, and M. Sinkinson (2021), ``Common Ownership in America: 1980--2017,'' \emph{American Economic Journal: Microeconomics} 13(3), 273--308. \href{https://www.aeaweb.org/articles?id=10.1257/mic.20190389}{Official article and abstract}.
\end{thebibliography}

\end{document}
